%==============================================================================
%  CHAPTER: System Design and Models
%  A diagram-rich chapter that visualises the architecture, behaviour and data
%  of the platform. Figures are generated TikZ/UML diagrams stored in
%  chapters/figs/ and pulled in with \input.
%==============================================================================
\chapter{System Design and Models}
\label{ch:design}

The previous chapter described \emph{what} the platform must do and the
principles that shape it. This chapter makes the design concrete by modelling
the system from six complementary viewpoints: the static \textbf{architecture}
of the deployed services, the \textbf{use cases} each role performs, the
\textbf{class} (domain) model of every subsystem, the
\textbf{entity-relationship} model of each database, the \textbf{behavioural
sequences} that exercise the architecture at run time, and the \textbf{state
machines} that govern the most important life cycles. Together these models
provide a complete blueprint of the Multi-Sport Club Management System (MSCMS)
and form the bridge between the requirements of
Chapter~\ref{ch:methodology} and the delivered system of
Chapter~\ref{ch:results}. All diagrams in this chapter are drawn to reflect the
system as actually implemented.

%==============================================================================
\section{System Architecture}
\label{sec:design-arch}

The platform is a cloud-native, microservices system in which a single Next.js
front-end communicates exclusively with a Spring Cloud API gateway, and the
gateway routes authenticated traffic to a set of independent domain services.
Figure~\ref{fig-arch-overview} presents the high-level architecture: the gateway
verifies the Keycloak-issued JWT on every request and forwards it to one of the
six domain services, each of which owns a private PostgreSQL database, so that no
service ever reaches into another's data. A Eureka/Config control plane provides
discovery and centralised configuration, Apache~Kafka carries domain events
asynchronously to the notification service, and a decoupled Python/FastAPI
machine-learning service and the external football-data.org feed sit at the
edges of the system. The raw chat WebSocket is the single channel that
deliberately bypasses the gateway.

% auto-generated figure: arch-overview (source: wf)
\begin{figure}[H]
\centering
\resizebox{\textwidth}{!}{%
\begin{tikzpicture}[node distance=6mm and 7mm,
    every node/.style={transform shape}]

  % ---- frontend (browser) ----
  \node[comp, minimum width=50mm] (fe) {Next.js / React Front-end (browser)};

  % ---- API gateway ----
  \node[svc, below=12mm of fe, minimum width=64mm] (gw)
        {Spring Cloud API Gateway\\(:8080 -- JWT verify + routing)};

  % ---- identity provider (left of gateway) ----
  \node[comp, left=18mm of gw, minimum width=28mm] (kc)
        {Keycloak\\realm \texttt{mscms}};

  % ---- control plane (right of gateway) ----
  \node[comp, right=18mm of gw, minimum width=28mm] (eu) {Eureka\\discovery (:8761)};
  \node[comp, below=6mm of eu, minimum width=28mm] (cfg) {Config\\Server (:8082)};

  % ---- six domain services in a row ----
  \node[svc, below=30mm of gw, xshift=-72mm, minimum width=24mm] (um)
        {user-\\management\\:8081};
  \node[svc, right=of um, minimum width=24mm] (pm) {player-\\management\\:8084};
  \node[svc, right=of pm, minimum width=24mm] (tm) {training-\\match\\:8087};
  \node[svc, right=of tm, minimum width=24mm] (mf) {medical-\\fitness\\:8086};
  \node[svc, right=of mf, minimum width=24mm] (ra) {reports-\\analytics\\:8083};
  \node[svc, right=of ra, minimum width=24mm] (nm) {notification-\\mail\\:8085};

  % ---- per-service databases ----
  \node[store, below=9mm of um] (dbu) {mscms\\\_user};
  \node[store, below=9mm of pm] (dbp) {mscms\\\_player};
  \node[store, below=9mm of tm] (dbt) {mscms\\\_training};
  \node[store, below=9mm of mf] (dbm) {mscms\\\_medical};
  \node[store, below=9mm of ra] (dbr) {mscms\\\_reports};
  \node[store, below=9mm of nm] (dbn) {mscms\\\_notif.};

  % ---- external football API (far top-left corner) ----
  \node[ext, left=14mm of fe, yshift=-2mm, minimum width=32mm] (fd)
        {football-data.org\\REST API (external)};

  % ---- async / infra row (bottom) ----
  \node[comp, below=12mm of dbt, minimum width=32mm] (kafka)
        {Apache Kafka\\(async event backbone)};
  \node[comp, left=10mm of kafka, minimum width=18mm] (redis) {Redis\\(cache)};
  \node[comp, below=12mm of dbn, minimum width=20mm] (rmq) {RabbitMQ};
  \node[comp, below=12mm of dbr, minimum width=28mm] (ml)
        {Python / FastAPI\\ML service};

  % ===== synchronous edges =====
  \draw[flow] (fe) -- node[lbl,pos=0.5]{HTTPS + JWT} (gw);
  \draw[flow] (gw.south) -- (um.north);
  \draw[flow] (gw.south) -- (pm.north);
  \draw[flow] (gw.south) -- (tm.north);
  \draw[flow] (gw.south) -- (mf.north);
  \draw[flow] (gw.south) -- (ra.north);
  \draw[flow] (gw.south) -- (nm.north);

  % service -> own db
  \draw[flow] (um) -- (dbu);
  \draw[flow] (pm) -- (dbp);
  \draw[flow] (tm) -- (dbt);
  \draw[flow] (mf) -- (dbm);
  \draw[flow] (ra) -- (dbr);
  \draw[flow] (nm) -- (dbn);

  % keycloak JWT to gateway + Admin API to user-management
  \draw[biflow] (kc) -- node[lbl,pos=0.5]{JWKS} (gw);
  \draw[flow] (kc.south) -- node[lbl,pos=0.5]{Admin API} (um.north -| kc.south) -- (um.north);

  % training-match <- external API (scheduled ingestion)
  \draw[flow] (fd.south) to[out=-90,in=150]
        node[lbl,pos=0.7]{scheduled ingest} (tm.north west);

  % reports -> ML
  \draw[flow] (ra) -- (ml);

  % notification raw websocket to frontend (outside gateway)
  \draw[flow] (nm.east) to[out=0,in=0,looseness=1.3]
        node[lbl,pos=0.5]{\texttt{/ws/chat} (raw WebSocket, outside gateway)} (fe.east);

  % ===== control-plane (dashed: register / pull config) =====
  \draw[flowd] (gw.east) -- node[lbl,pos=0.5]{register} (eu.west);
  \draw[flowd] (cfg.west) to[out=180,in=-45] node[lbl,pos=0.7]{config} (gw.south east);
  \foreach \s in {pm,tm,mf,ra,nm}{ \draw[flowd] (\s.north) to[out=70,in=-130] (eu.south west); }

  % ===== async Kafka backbone =====
  \draw[flowd] (tm.south) to[out=-90,in=90] (kafka.north);
  \draw[flowd] (mf.south) to[out=-90,in=60] (kafka.north east);
  \draw[flowd] (pm.south) to[out=-90,in=150] (kafka.north west);
  \draw[flowd] (kafka.east) to[out=0,in=-135] node[lbl,pos=0.6]{consume} (nm.south);

\end{tikzpicture}%
}
\caption{High-level cloud-native microservices architecture of the platform: a Next.js front-end talks only to the Spring Cloud API Gateway, which routes authenticated traffic to six domain services, each owning a private PostgreSQL database, supported by a Eureka/Config control plane, Keycloak identity, a Kafka event backbone, Redis/RabbitMQ infrastructure, a Python/FastAPI ML service, and the external football-data.org feed.}
\label{fig-arch-overview}
\end{figure}


Figure~\ref{fig-arch-deploy} shows the same system from a deployment standpoint:
every component is a container orchestrated by Docker Compose on a single shared
network, grouped into infrastructure backing services, the Spring Cloud control
plane, the six domain services, the machine-learning service and the front-end.
This is the unit of reproducible deployment described in Appendix~\ref{app:deploy}.

\input{chapters/figs/arch-deploy}

%==============================================================================
\section{Use-Case Models}
\label{sec:design-usecase}

The platform serves many roles, each with a distinct surface of operations.
The following use-case diagrams group those operations by the actors that
perform them. Figure~\ref{fig-uc-admin} covers administration and club
management (the administrator, sport manager and team manager);
Figure~\ref{fig-uc-match} covers match and training operations (coaching staff
and analysts); Figure~\ref{fig-uc-medical} covers the medical and fitness staff;
and Figure~\ref{fig-uc-business} covers scouting, sponsorship and the read-only
surface presented to players and fans.

\input{chapters/figs/uc-admin}
% auto-generated figure: uc-match (source: wf)
\begin{figure}[H]
\centering
\resizebox{\textwidth}{!}{%
\begin{tikzpicture}
  % actors -- Coach on the left, Analyst on the right
  \UCactor{0}{0}{Coach}
  \UCactor{20}{0}{Analyst}

  % use cases -- training & match operations (Coach)
  \node[usecase] (u1) at (6.5,4.0){Create training plan,\\ session \& drills};
  \node[usecase] (u2) at (6.5,2.4){Record attendance};
  \node[usecase] (u3) at (6.5,0.8){Capture player\\ assessments};
  \node[usecase] (u4) at (6.5,-0.8){Schedule official match\\ from a fixture};
  \node[usecase] (u5) at (6.5,-2.4){Build starting line-up\\ (formation)};
  \node[usecase] (u6) at (6.5,-4.0){Run live match (goals,\\ cards, subs, injuries)};

  % use cases -- analysis & ML (Analyst)
  \node[usecase] (u7) at (13.5,0.8){Author match \&\\ player analysis};
  \node[usecase] (u8) at (13.5,-0.8){Consume ML outcome \&\\ player-rating predictions};

  % Coach associations
  \draw[dLine] (0.5,-0.4) -- (u1);
  \draw[dLine] (0.5,-0.4) -- (u2);
  \draw[dLine] (0.5,-0.4) -- (u3);
  \draw[dLine] (0.5,-0.4) -- (u4);
  \draw[dLine] (0.5,-0.4) -- (u5);
  \draw[dLine] (0.5,-0.4) -- (u6);

  % Analyst associations
  \draw[dLine] (19.5,-0.4) -- (u7);
  \draw[dLine] (19.5,-0.4) -- (u8);
\end{tikzpicture}%
}
\caption{Use-case diagram for match and training operations, in which the coach drives training planning and the line-up-mandatory match and live-match workflows while the analyst authors match and player analyses and consumes the machine-learning match-outcome and player-rating predictions.}
\label{fig-uc-match}
\end{figure}

% auto-generated figure: uc-medical (source: wf)
\begin{figure}[H]
\centering
\resizebox{\textwidth}{!}{%
\begin{tikzpicture}
  % --- actors (left column) ---
  \UCactor{0}{4.4}{Doctor}
  \UCactor{0}{0}{Physiotherapist}
  \UCactor{0}{-4.4}{Fitness Coach}

  % --- use cases (right column) ---
  \node[usecase] (u1) at (6,6.0){Log injury};
  \node[usecase] (u2) at (6,4.4){Diagnose injury};
  \node[usecase] (u3) at (6,2.8){Prescribe treatment};
  \node[usecase] (u4) at (6,1.0){Manage rehabilitation};
  \node[usecase] (u5) at (6,-0.6){Manage recovery programme};
  \node[usecase] (u6) at (6,-3.4){Run fitness test};
  \node[usecase] (u7) at (6,-5.0){View Treatment Room};
  \node[usecase] (u8) at (10.4,0.4){Maintain medical\\ catalogues};

  % --- Doctor -> log / diagnose / treat ---
  \draw[dLine] (0.5,4.0) -- (u1);
  \draw[dLine] (0.5,4.0) -- (u2);
  \draw[dLine] (0.5,4.0) -- (u3);

  % --- Physiotherapist -> rehabilitation / recovery ---
  \draw[dLine] (0.5,-0.4) -- (u4);
  \draw[dLine] (0.5,-0.4) -- (u5);

  % --- Fitness Coach -> fitness test + treatment room ---
  \draw[dLine] (0.5,-4.8) -- (u6);
  \draw[dLine] (0.5,-4.8) -- (u7);

  % --- shared catalogue maintenance ---
  \draw[dLine] (u3) -- (u8);
  \draw[dLine] (u5) -- (u8);
  \draw[dLine] (u6) -- (u8);
\end{tikzpicture}%
}
\caption{Use-case diagram for the medical and fitness module, mapping the Doctor, Physiotherapist and Fitness Coach to the staged injury-to-recovery workflow and the shared medical catalogues.}
\label{fig-uc-medical}
\end{figure}

\input{chapters/figs/uc-business}

%==============================================================================
\section{Class and Domain Models}
\label{sec:design-class}

The relational schema is materialised in the running services as a set of JPA
entity classes related by composition (an owning parent that manages the life
cycle of its children) and by reference (an identifier the application resolves
on demand). The class diagrams below capture each subsystem in turn. The user
identity model (Figure~\ref{fig-class-user}) shows the central \texttt{User}
specialised by a one-to-one profile per role. The squad model
(Figure~\ref{fig-class-squad}) expresses the team hierarchy and the transfer
market. The match and training model (Figure~\ref{fig-class-match}) and the
competition model (Figure~\ref{fig-class-competition}) show why a match, its
formation and its starting line-up form a single atomic aggregate. The medical
model (Figure~\ref{fig-class-medical}) encodes the recovery journey as a
composition hierarchy rooted at \texttt{Injury}, and the communication
(Figure~\ref{fig-class-comm}) and analytics (Figure~\ref{fig-class-reports})
models complete the picture.

% auto-generated figure: class-user (source: wf)
\begin{figure}[H]
\centering
\resizebox{\textwidth}{!}{%
\begin{tikzpicture}[node distance=10mm]
  % --- central identity class ---
  \UmlClass{user}{User}{%
    - keycloakId : String\\
    - username : String\\
    - email : String\\
    - role : Role\\
    - age : int\\
    - gender : Gender\\
    - phone : String}{%
    + assignRole()\\
    + profile()}

  % --- role-extension classes (vertical extension, 1..1) ---
  \UmlClass[above right=2mm and 30mm of user]{player}{PlayerProfile}{%
    - dateOfBirth : Date\\
    - nationality : String\\
    - preferredPosition\\
    - marketValue : Money\\
    - kitNumber : int\\
    - status : PlayerStatus\\
    - rosterId : Long\\
    - contractId : Long}{}

  \UmlClass[right=30mm of user]{coach}{HeadCoachProfile}{%
    - sport : Sport\\
    - team : String\\
    - staffRole : String\\
    - experienceYears : int\\
    - licenceLevel : String}{}

  \UmlClass[below right=2mm and 30mm of user]{nteam}{NationalTeam}{%
    - federationName : String\\
    - country : String\\
    - contact : String}{}

  % --- note: further role extensions ---
  \UmlClass[left=22mm of user]{others}{Other extensions}{%
    Doctor, Physiotherapist,\\
    FitnessCoach, TeamManager,\\
    SportManager, Scout,\\
    Sponsor, Fan\\
    \emph{(same 1..1 pattern)}}{}

  % --- composition links: User extends each profile 1..1 ---
  \draw[flow] (user.east) -- node[lbl,pos=0.55]{\itshape extends 1..1} (player.west);
  \draw[flow] (user.east) -- node[lbl,pos=0.55]{\itshape extends 1..1} (coach.west);
  \draw[flow] (user.east) -- node[lbl,pos=0.55]{\itshape extends 1..1} (nteam.west);
  \draw[flow] (user.west) -- node[lbl,pos=0.5]{\itshape extends 1..1} (others.east);
\end{tikzpicture}}
\caption{Class model of the user-management domain: the central \texttt{User} identity record and its one-to-one role-extension classes (\texttt{PlayerProfile}, \texttt{HeadCoachProfile}, \texttt{NationalTeam} and the remaining role extensions), realising the vertical-extension pattern of \texttt{mscms\_user}.}
\label{fig-class-user}
\end{figure}

\input{chapters/figs/class-squad}
\input{chapters/figs/class-match}
% auto-generated figure: class-competition (source: wf)
\begin{figure}[H]
\centering
\begin{tikzpicture}[node distance=10mm]
  \UmlClass{comp}{Competition}{%
    - name : String\\
    - type : CompetitionType\\
    - status : CompetitionStatus}{%
    + isCompleted()\\
    + standings()}
  \UmlClass[below left=16mm and -22mm of comp]{team}{CompetitionTeam}{%
    - name : String\\
    - shortName : String\\
    - crestUrl : String\\
    - sport : Sport}{}
  \UmlClass[below right=16mm and -22mm of comp]{fix}{CompetitionFixture}{%
    - homeTeamId : Long\\
    - awayTeamId : Long\\
    - group : String [0..1]\\
    - matchday : int\\
    - round : String\\
    - homeScore : Integer\\
    - awayScore : Integer}{%
    + isPlayed()}

  % compositions (filled diamond at the Competition / whole end)
  \draw[{Diamond[fill=dLine,length=3mm,width=2mm]}-,draw=dLine,thick]
        (comp.south) -- node[lbl,pos=0.85,left]{1..*} (team.north);
  \draw[{Diamond[fill=dLine,length=3mm,width=2mm]}-,draw=dLine,thick]
        (comp.south) -- node[lbl,pos=0.85,right]{1..* \{ordered\}} (fix.north);

  \node[lbl,align=left,text width=52mm,font=\scriptsize,
        below=4mm of fix.south,anchor=north]
       {\itshape round distinguishes the Champions League league
        phase (e.g.\ ``League Phase'') from the knockout rounds
        (e.g.\ ``Round of 16'', ``Final'').};
\end{tikzpicture}
\caption{Class model of the competitions domain: a Competition composes its participating CompetitionTeam rows and its ordered CompetitionFixture rows, whose round label separates the Champions League league phase from the knockout rounds.}
\label{fig-class-competition}
\end{figure}

\input{chapters/figs/class-medical}
% auto-generated figure: class-comm (source: wf)
\begin{figure}[H]
\centering
\begin{tikzpicture}[node distance=10mm]
  \UmlClass{notif}{Notification}{%
    - recipientKeycloakId\\
    - type\\
    - category\\
    - title\\
    - body\\
    - read : boolean\\
    - createdAt}{}
  \UmlClass[right=10mm of notif]{alert}{Alert}{%
    - type\\
    - priority\\
    - targetRole : opt.\\
    - acknowledged : boolean\\
    - resolved : boolean}{}
  \UmlClass[below=12mm of notif]{msg}{Message}{%
    - senderKeycloakId\\
    - groupId\\
    - body\\
    - parentMessageId : opt.\\
    - createdAt}{}
  \UmlClass[right=22mm of msg]{grp}{ChatGroup}{%
    - name\\
    - ownerKeycloakId}{}
  \UmlClass[below=12mm of grp]{mem}{ChatGroupMember}{%
    - memberKeycloakId\\
    - status : invited / member\\
    - invitedBy}{}

  % composition: ChatGroup composes 1..* ChatGroupMember
  \draw[flow] (grp) -- node[lbl,pos=0.5]{\itshape composes 1..*} (mem);
  % reference: Message -> ChatGroup
  \draw[flow] (msg) -- node[lbl,pos=0.5]{\itshape posted in 1} (grp.west);
\end{tikzpicture}
\caption{Class model of the communication subsystem (\texttt{mscms\_notification}), showing per-recipient notifications, role-targeted alerts, threaded messages, and chat groups that compose their members.}
\label{fig-class-comm}
\end{figure}

% auto-generated figure: class-reports (source: wf)
\begin{figure}[H]
\centering
\resizebox{\textwidth}{!}{%
\begin{tikzpicture}[node distance=10mm]
  \UmlClass{ma}{MatchAnalysis}{%
    - matchId : Long\\
    - teamId : Long\\
    - sportType : SportType\\
    - sportSpecificStats : json\\
    - tacticalAnalysis : json\\
    - playerRatings : json}{}
  \UmlClass[right=12mm of ma]{pa}{PlayerAnalytics}{%
    - playerKeycloakId\\
    - teamId : Long\\
    - sportType : SportType\\
    - primaryScore : double\\
    - secondaryScore : double\\
    - averageRating : double}{}
  \UmlClass[right=12mm of pa]{ta}{TeamAnalytics}{%
    - teamId : Long\\
    - sportType : SportType\\
    - wins / draws / losses\\
    - pointsFor : int\\
    - pointsAgainst : int}{}
  \UmlClass[below=14mm of ma]{tr}{TrainingAnalytics}{%
    - teamId : Long\\
    - period : DateRange\\
    - attendanceRate\\
    - completionRate\\
    - trendData : json}{}
  \UmlClass[below=14mm of pa]{sr}{ScoutReport}{%
    - scoutKeycloakId\\
    - outerPlayerId\\
    - ratings : json\\
    - recommendation : bool\\
    - notes : text}{}
  \UmlClass[below=14mm of ta]{so}{SponsorContractOffer}{%
    - sponsorKeycloakId\\
    - teamId : Long\\
    - contractValue\\
    - durationMonths : int\\
    - terms : text\\
    - status : OfferStatus}{}

  % relationship notes: opaque-id references, no object links
  \draw[flowd] (ma) -- node[lbl,pos=0.5]{\itshape matchId} (tr);
  \draw[flowd] (pa) -- node[lbl,pos=0.5]{\itshape playerKeycloakId} (sr);
  \node[lbl,below=8mm of sr,text width=70mm,align=center]
        {\itshape all classes reference operational entities
         (matches, players, teams, sponsors) by opaque id --- no object links};
\end{tikzpicture}}
\caption{Class model of the reports-and-analytics domain (\texttt{mscms\_reports}): six authored or derived classes that aggregate operational data over a period and reference matches, players, teams and sponsors only by opaque identifier rather than by object links.}
\label{fig-class-reports}
\end{figure}


%==============================================================================
\section{Entity-Relationship (Database) Models}
\label{sec:design-er}

Because every domain service owns its own database, the data model is best
presented one database at a time. Each diagram shows the principal tables of a
service database, their primary and foreign keys and the relationships between
them. Figure~\ref{fig-er-user} models the user database with its central
\texttt{user} table and per-role extension tables;
Figure~\ref{fig-er-player} the squad and transfer market;
Figure~\ref{fig-er-training} the (largest) training, match and competition
schema; Figure~\ref{fig-er-medical} the injury aggregate and independent
measurement tables; Figure~\ref{fig-er-notification} the notification,
alerting and group-chat tables; and Figure~\ref{fig-er-reports} the derived
analytics and authored artefacts, which reference the operational domains only
by opaque identifier.

% auto-generated figure: er-user (source: wf)
\begin{figure}[H]
\centering
\resizebox{\textwidth}{!}{%
\begin{tikzpicture}[node distance=14mm]
  \Entity{usr}{USER}{%
    \underline{PK} id\\ keycloakId\\ username\\ email\\ role : Role\\ age\\ gender\\ phone}
  \Entity[right=34mm of usr]{ply}{PLAYER}{%
    \underline{PK} id\\ \underline{FK} userId\\ dateOfBirth\\ nationality\\ position\\ marketValue\\ kitNumber\\ status}
  \Entity[below left=10mm and 6mm of usr]{coa}{HEAD\_COACH}{%
    \underline{PK} id\\ \underline{FK} userId\\ sport\\ team\\ staffRole\\ licenceLevel}
  \Entity[below right=10mm and 6mm of usr]{nat}{NATIONAL\_TEAM}{%
    \underline{PK} id\\ federationName\\ country\\ contact}
  \Rel{usr}{ply}{1}{1}{extends}
  \Rel{usr}{coa}{1}{1}{extends}
  \Rel{usr}{nat}{1}{1}{extends}
  \node[font=\scriptsize\itshape, text width=46mm, align=left,
        below=22mm of usr] (note)
    {Further 1..1 role extensions follow the same pattern: doctor, physiotherapist,
     fitness\_coach, team\_manager, sport\_manager, scout, sponsor and fan.};
\end{tikzpicture}%
}
\caption{Entity-relationship diagram of the \texttt{mscms\_user} database, showing the common USER identity record and its one-to-one per-role profile extensions (PLAYER, HEAD\_COACH, NATIONAL\_TEAM, and the remaining role tables).}
\label{fig-er-user}
\end{figure}

\input{chapters/figs/er-player}
% auto-generated figure: er-training (source: wf)
\begin{figure}[H]
\centering
\resizebox{\textwidth}{!}{%
\begin{tikzpicture}[node distance=10mm,
    every node/.append style={font=\scriptsize}]

  % ===== Matches cluster (left/centre) ==================================
  \Entity{m}{MATCHES}{%
    \underline{PK} id\\
    home\_team\_id\\
    opponent\_name\\
    status, sport\_type\\
    \underline{FK} competition\_id\\
    home/away\_score\\
    kickoff\_time}
  % formations / lineups stacked to the right of MATCHES
  \Entity[right=34mm of m]{mf}{MATCH\_FORMATIONS}{%
    \underline{PK} id\\
    team\_id\\
    formation\\
    tactical\_approach}
  \Entity[below=14mm of mf]{ml}{MATCH\_LINEUPS}{%
    \underline{PK} id\\
    \underline{FK} match\_id\\
    player\_keycloak\_id\\
    position, starter}
  % events / statistics stacked below MATCHES
  \Entity[below=16mm of m]{me}{MATCH\_EVENTS}{%
    \underline{PK} id\\
    \underline{FK} match\_id\\
    event\_type\\
    minute, extra\_time}
  \Entity[below=14mm of me]{pms}{PLAYER\_MATCH\_\\STATISTICS}{%
    \underline{PK} id\\
    \underline{FK} match\_id\\
    player\_keycloak\_id\\
    sport\_type, stats}

  % ===== Competition cluster (right) ====================================
  \Entity[right=40mm of mf]{cmp}{COMPETITION}{%
    \underline{PK} id\\
    name, season\\
    type, rounds}
  \Entity[below=16mm of cmp]{ct}{COMPETITION\_TEAM}{%
    \underline{PK} id\\
    \underline{FK} competition\_id\\
    name, sport\\
    crest\_url}
  \Entity[below=14mm of ct]{cf}{COMPETITION\_FIXTURE}{%
    \underline{PK} id\\
    \underline{FK} competition\_id\\
    home/away\_score\\
    round, matchday, played}

  % ===== Training cluster (bottom row) ==================================
  \Entity[below=24mm of pms]{ts}{TRAINING\_SESSIONS}{%
    \underline{PK} id\\
    team\_id\\
    session\_type, status\\
    scheduled\_date}
  \Entity[right=34mm of ts]{tp}{TRAINING\_PLANS}{%
    \underline{PK} id\\
    \underline{FK} training\_session\_id\\
    category, intensity}
  \Entity[below=14mm of tp]{td}{TRAINING\_DRILLS}{%
    \underline{PK} id\\
    \underline{FK} training\_session\_id\\
    name, duration}
  \Entity[right=34mm of tp]{ta}{TRAINING\_ATTENDANCE}{%
    \underline{PK} id\\
    \underline{FK} training\_session\_id\\
    player\_keycloak\_id\\
    attended, perf\_score}

  % ===== Relationships : MATCHES hub ===================================
  % to formations (right) and lineups (lower-right) from the right edge
  \draw[dLine] (m.east) -- node[lbl,pos=0.18]{1} node[lbl,pos=0.55]{\itshape set by}
       node[lbl,pos=0.9]{1} (mf.west);
  \draw[dLine] (m.east) -- node[lbl,pos=0.12]{1} node[lbl,pos=0.6]{\itshape names}
       node[lbl,pos=0.92]{*} (ml.west);
  % to events / statistics straight down
  \Rel{m}{me}{1}{*}{records}
  \draw[dLine] (me.south) -- node[lbl,pos=0.5]{\itshape measures}
       node[lbl,pos=0.9]{*} (pms.north);

  % ===== Relationships : COMPETITION hub ===============================
  % COMPETITION 1--* MATCHES (competition_id FK) routed over the top so it
  % never crosses the formations/lineups boxes
  \draw[dLine] (cmp.north) |- ($(m.north)+(0,7mm)$)
       node[lbl,pos=0.55]{\itshape groups} -- (m.north);
  \node[lbl] at ($(cmp.north)+(0,4mm)$){1};
  \node[lbl] at ($(m.north)+(3mm,4mm)$){*};
  % vertical chain, offset so it never crosses a box
  \draw[dLine] (cmp.south) -- node[lbl,pos=0.2]{1} node[lbl,pos=0.55]{\itshape includes}
       node[lbl,pos=0.85]{*} (ct.north);
  \draw[dLine] (ct.east) -- ++(8mm,0) |- node[lbl,pos=0.75]{\itshape schedules} (cf.east);
  \node[lbl] at ($(ct.east)+(10mm,-7mm)$){*};

  % ===== Relationships : TRAINING hub ==================================
  \draw[dLine] (ts.east) -- node[lbl,pos=0.18]{1} node[lbl,pos=0.55]{\itshape follows}
       node[lbl,pos=0.9]{*} (tp.west);
  \draw[dLine] (tp.south) -- node[lbl,pos=0.5]{\itshape runs}
       node[lbl,pos=0.85]{*} (td.north);
  \draw[dLine] (tp.east) -- node[lbl,pos=0.5]{\itshape tracks}
       node[lbl,pos=0.9]{*} (ta.west);

\end{tikzpicture}%
}
\caption{Entity-relationship diagram of the \texttt{mscms\_training} database, showing the matches subsystem (formations, line-ups, events and per-player statistics), the database-backed competition subsystem (teams and fixtures), and the training subsystem (plans, drills and attendance) with their principal columns and foreign-key relationships.}
\label{fig-er-training}
\end{figure}

% auto-generated figure: er-medical (source: wf)
\begin{figure}[H]
\centering
\resizebox{\textwidth}{!}{%
\begin{tikzpicture}[node distance=14mm]
  % central injury record and its staged journey
  \Entity{inj}{INJURIES}{%
    \underline{PK} id\\
    playerKeycloakId\\
    teamId\\
    injuryType\\
    severity\\
    injuredBodyPart\\
    injuryDate\\
    expectedRecoveryDate\\
    status}
  \Entity[right=30mm of inj]{dia}{DIAGNOSES}{%
    \underline{PK} id\\
    \underline{FK} injuryId\\
    playerKeycloakId\\
    diagnosis\\
    testResults\\
    recommendations}
  \Entity[above=12mm of dia]{trt}{TREATMENTS}{%
    \underline{PK} id\\
    \underline{FK} injuryId\\
    treatmentType\\
    status\\
    startDate\\
    endDate}
  \Entity[below=12mm of dia]{reh}{REHABILITATIONS}{%
    \underline{PK} id\\
    \underline{FK} injuryId\\
    playerKeycloakId\\
    programmeDetail\\
    status}
  \Entity[right=30mm of reh]{rec}{RECOVERY\_PROGRAMS}{%
    \underline{PK} id\\
    \underline{FK} rehabilitationId\\
    playerKeycloakId\\
    programmeDetail}
  % independent fitness / load entities
  \Entity[below=14mm of reh]{fit}{FITNESS\_TESTS}{%
    \underline{PK} id\\
    playerKeycloakId\\
    teamId\\
    testType\\
    sportType\\
    testDate\\
    result\\
    resultCategory}
  \Entity[right=30mm of fit]{tld}{TRAINING\_LOADS}{%
    \underline{PK} id\\
    playerKeycloakId\\
    intensity\\
    load}
  % relationships
  \Rel{inj}{trt}{1}{*}{has}
  \Rel{inj}{dia}{1}{*}{has}
  \Rel{inj}{reh}{1}{*}{has}
  \Rel{reh}{rec}{1}{*}{drives}
\end{tikzpicture}%
}
\caption{Entity-relationship diagram for the \texttt{mscms\_medical} database, showing the staged injury journey (injuries owning diagnoses, treatments and rehabilitations, with rehabilitations driving recovery programmes) alongside the independent fitness-test and training-load records.}
\label{fig-er-medical}
\end{figure}

% auto-generated figure: er-notification (source: wf)
\begin{figure}[H]
\centering
\resizebox{\textwidth}{!}{%
\begin{tikzpicture}[node distance=16mm]
  \Entity{cg}{CHAT\_GROUP}{%
    \underline{PK} id\\ name\\ \underline{FK} ownerKeycloakId}
  \Entity[below=18mm of cg]{cgm}{CHAT\_GROUP\_MEMBER}{%
    \underline{PK} id\\ \underline{FK} groupId\\ memberKeycloakId\\ status\\ invitedBy}
  \Entity[right=34mm of cg]{msg}{MESSAGES}{%
    \underline{PK} id\\ senderKeycloakId\\ \underline{FK} groupId\\ body\\ \underline{FK} parentMessageId}
  \Entity[right=34mm of cgm]{ntf}{NOTIFICATIONS}{%
    \underline{PK} id\\ recipientKeycloakId\\ type\\ category\\ title\\ body\\ read}
  \Entity[below=18mm of cgm]{alr}{ALERTS}{%
    \underline{PK} id\\ type\\ priority\\ targetRole\\ acknowledged\\ resolved}
  \Rel{cg}{cgm}{1}{*}{has}
  \Rel{cg}{msg}{1}{*}{contains}
  \draw[dLine] (msg.north east) .. controls +(1.4,0.9) and +(1.4,-0.2) ..
       node[lbl,pos=0.5,xshift=4mm]{\itshape replies to 0..1} (msg.east);
\end{tikzpicture}%
}
\caption{Entity-relationship diagram of the \texttt{mscms\_notification} database, showing notifications, alerts, and the team-chat groups, members, and threaded messages.}
\label{fig-er-notification}
\end{figure}

% auto-generated figure: er-reports (source: wf)
\begin{figure}[H]
\centering
\begin{tikzpicture}[node distance=8mm]
  \Entity{ma}{MATCH\_ANALYSIS}{%
    \underline{PK} id\\
    matchId\\
    teamId\\
    sportType\\
    sportSpecificStats\\
    tacticalAnalysis\\
    playerRatings}
  \Entity[right=10mm of ma]{pa}{PLAYER\_ANALYTICS}{%
    \underline{PK} id\\
    playerKeycloakId\\
    teamId\\
    sportType\\
    primaryScore\\
    secondaryScore\\
    averageRating}
  \Entity[right=10mm of pa]{ta}{TEAM\_ANALYTICS}{%
    \underline{PK} id\\
    teamId\\
    sportType\\
    wins / draws / losses\\
    pointsFor\\
    pointsAgainst}
  \Entity[below=8mm of ma]{tr}{TRAINING\_ANALYTICS}{%
    \underline{PK} id\\
    teamId\\
    aggregate training\\ \quad metrics}
  \Entity[below=8mm of pa]{sr}{SCOUT\_REPORTS}{%
    \underline{PK} id\\
    scoutKeycloakId\\
    outerPlayerId\\
    ratings\\
    recommendation}
  \Entity[below=8mm of ta]{sc}{SPONSOR\_CONTRACT\_OFFERS}{%
    \underline{PK} id\\
    sponsorKeycloakId\\
    offerTitle\\
    contractValue\\
    sponsorshipType\\
    status}
  \node[lbl, draw=dLine, fill=dFill, text width=46mm, align=center,
        below=11mm of sr] (note)
    {\itshape Entities are independent; cross-domain links
     (matchId, teamId, player\slash scout\slash sponsor Keycloak ids)
     are stored as opaque references, not FK constraints.};
\end{tikzpicture}
\caption{Entity-relationship diagram of the \texttt{mscms\_reports} database, whose analytical entities aggregate operational data from other domains through opaque cross-service references rather than foreign-key constraints.}
\label{fig-er-reports}
\end{figure}


%==============================================================================
\section{Behavioural Sequence Models}
\label{sec:design-seq}

The static models above are exercised at run time by a number of behavioural
flows. The sequence diagrams in this section trace the cooperation between the
front-end, the gateway, the domain services, Keycloak and Kafka for the most
important interactions. They cover authentication and token issuance
(Figure~\ref{fig-seq-login}); the atomic scheduling of an official match with
its line-up (Figure~\ref{fig-seq-schedule}); live-match phase progression and
the minute-bounded goal event (Figure~\ref{fig-seq-live}); an in-match injury
that forces a substitution and creates a real medical record across two services
(Figure~\ref{fig-seq-injury-sub}); the staged injury-to-recovery workflow
(Figure~\ref{fig-seq-medical}); real-time group messaging over WebSocket with a
polling fallback (Figure~\ref{fig-seq-messaging}); the scheduled ingestion of
external fixtures (Figure~\ref{fig-seq-ingest}); event-driven notification
creation over Kafka (Figure~\ref{fig-seq-notify}); and a password change that is
only permitted after the current password is verified
(Figure~\ref{fig-seq-password}).

% auto-generated figure: seq-login (source: wf)
\begin{figure}[H]
\centering
{\footnotesize
\begin{sequencediagram}
  \newthread{u}{:User}
  \newinst[1]{fe}{:Frontend}
  \newinst[1]{gw}{:Gateway}
  \newinst[1]{um}{:user-management}
  \newinst[1]{kc}{:Keycloak}
  \begin{call}{u}{submit username/password}{fe}{token + role stored}
    \begin{call}{fe}{POST /auth/login}{gw}{access token, role}
      \begin{call}{gw}{forward login}{um}{access token, role}
        \begin{call}{um}{OAuth2 password-grant (token endpoint)}{kc}{access + refresh tokens}
        \end{call}
      \end{call}
    \end{call}
  \end{call}
  \postlevel
  \mess{fe}{decode JWT: sub=keycloakId, realm roles}{fe}
\end{sequencediagram}
}
\caption{Sequence diagram of the authentication flow: credentials submitted on the front-end traverse the gateway and the user-management service, which performs an OAuth~2.0 password-grant exchange at Keycloak to obtain and return the JWT access token and role.}
\label{fig-seq-login}
\end{figure}

% auto-generated figure: seq-schedule (source: wf)
\begin{figure}[H]
\centering
{\footnotesize
\begin{sequencediagram}
  \newthread{u}{:Coach}
  \newinst[1]{fe}{:Frontend}
  \newinst[1]{gw}{:Gateway}
  \newinst[1]{svc}{:training-match}
  \newinst[1]{db}{:DB}

  \begin{call}{u}{pick fixture, build formation + starting XI}{fe}{scheduled match shown}
    \begin{call}{fe}{POST /matches (+formation, lineup)}{gw}{201 Created}
      \begin{call}{gw}{route createMatch()}{svc}{Match}
        \begin{sdblock}{validate}{kickoff $\geq$ 5\,min ahead, lineup present}
          \begin{messcall}{svc}{[invalid] 400 Bad Request}{svc}{}
          \end{messcall}
          \begin{call}{svc}{[ok] persist atomically (one tx)}{db}{Match + MatchFormation + MatchLineup}
          \end{call}
        \end{sdblock}
      \end{call}
    \end{call}
  \end{call}
\end{sequencediagram}
}
\caption{Sequence diagram for scheduling an official match: the coach builds the formation and starting eleven from a competition fixture, the training-match service validates the kick-off time and line-up, and on success persists the Match, MatchFormation and MatchLineup atomically in a single transaction so a match can never exist without a starting eleven.}
\label{fig-seq-schedule}
\end{figure}

% auto-generated figure: seq-live (source: wf)
\begin{figure}[H]
\centering
{\footnotesize
\begin{sequencediagram}
  \newthread{op}{:Operator}
  \newinst[1]{fe}{:Frontend}
  \newinst[1]{gw}{:Gateway}
  \newinst[1]{svc}{:training-match}
  \newinst[1]{db}{:DB}

  \begin{call}{op}{startMatch()}{fe}{phase: FirstHalf}
  \end{call}

  \begin{sdblock}{Phase machine}{FirstHalf $\to$ HalfTime $\to$ SecondHalf $\to$ FullTime}
    \begin{call}{op}{advancePhase()}{fe}{phase = next}
    \end{call}
  \end{sdblock}

  \begin{sdblock}{Record goal}{POST /match-events \{type=GOAL, minute\}}
    \begin{call}{op}{recordGoal(minute)}{fe}{updated state}
      \begin{call}{fe}{POST /match-events}{gw}{200 OK}
        \begin{call}{gw}{create(event)}{svc}{result}
          \begin{sdblock}{alt}{minute out of phase range}
            \begin{call}{svc}{reject: invalid minute}{svc}{422}
            \end{call}
          \end{sdblock}
          \begin{sdblock}{else}{minute within phase}
            \begin{call}{svc}{persist match\_event (minute, event\_type)}{db}{row}
            \end{call}
            \begin{call}{svc}{update home/away\_team\_score}{db}{ok}
            \end{call}
          \end{sdblock}
        \end{call}
      \end{call}
    \end{call}
  \end{sdblock}

\end{sequencediagram}
}
\caption{Sequence diagram of live-match phase progression and goal recording, showing the broadcast phase machine (first half, half-time, second half, full-time) and the minute-bounding validation that rejects out-of-phase goals before persisting the event and updating the score.}
\label{fig-seq-live}
\end{figure}

% auto-generated figure: seq-injury-sub (source: wf)
\begin{figure}[H]
\centering
{\scriptsize
\begin{sequencediagram}
  \newthread{op}{:Operator}
  \newinst[1]{fe}{:Frontend}
  \newinst[1]{gw}{:Gateway}
  \newinst[1]{tm}{:training-match}
  \newinst[1]{mf}{:medical-fitness}
  \newinst[1]{db}{:DB}
  \begin{call}{op}{record INJURY event(player, minute)}{fe}{substitution required}
    \begin{call}{fe}{POST /match-events}{gw}{201}
      \begin{call}{gw}{create(event\_type=INJURY)}{tm}{event}
        \begin{call}{tm}{persist match\_event}{db}{ok}
        \end{call}
        \begin{call}{tm}{POST /injuries(player, status=INJURED)}{mf}{injury}
          \begin{call}{mf}{persist injury (player flipped injured)}{db}{ok}
          \end{call}
        \end{call}
      \end{call}
    \end{call}
  \end{call}
  \begin{call}{op}{record SUBSTITUTION(out=injured, in=sub, minute)}{fe}{updated}
    \begin{call}{fe}{POST /match-events + PUT /match-lineups}{gw}{200}
      \begin{call}{gw}{updateLineup(playerOut@minute, subIn)}{tm}{lineup}
        \begin{call}{tm}{persist match\_event + match\_lineups}{db}{ok}
        \end{call}
      \end{call}
    \end{call}
  \end{call}
\end{sequencediagram}
}
\caption{Sequence diagram for an in-match injury and the forced substitution it triggers, showing how training-match persists the INJURY match event while delegating to medical-fitness to create the real injury record before the operator records the substitution.}
\label{fig-seq-injury-sub}
\end{figure}

\input{chapters/figs/seq-medical}
% auto-generated figure: seq-messaging (source: wf)
\begin{figure}[H]
\centering
{\footnotesize
\begin{sequencediagram}
  \newthread{u}{:User}
  \newinst[1]{fe}{:Frontend}
  \newinst[1.6]{ws}{:WSHandler /ws/chat}
  \newinst[1.4]{db}{:DB}
  \newinst[1.2]{mb}{:Members}

  \begin{call}{u}{open chat}{fe}{groups + history}
    \begin{call}{fe}{GET /messages/groups?user=}{ws}{groups}
    \end{call}
  \end{call}

  \mess{fe}{connect ws://.../ws/chat}{ws}
  \postlevel

  \begin{sdblock}{alt}{[\,WebSocket available\,]}
    \mess{u}{send content + groupId}{fe}
    \mess{fe}{message frame}{ws}
    \postlevel
    \begin{call}{ws}{persist Message(group\_id)}{db}{saved id}
    \end{call}
    \mess{ws}{broadcast to group}{mb}
    \postlevel
  \end{sdblock}

  \begin{sdblock}[dFill]{alt}{[\,socket unavailable\,]}
    \mess{u}{send content + groupId}{fe}
    \begin{call}{fe}{POST /messages \{group\_id\}}{ws}{saved}
      \begin{call}{ws}{persist Message}{db}{row}
      \end{call}
    \end{call}
    \begin{call}{fe}{poll GET /messages?group=}{ws}{messages}
    \end{call}
    \mess{fe}{render de-duplicated}{mb}
    \postlevel
  \end{sdblock}
\end{sequencediagram}
}
\caption{Sequence diagram for group messaging: the front-end opens a WebSocket to \texttt{/ws/chat} on the notification-mail service (outside the gateway) and, when the socket is available, sends a message that the handler persists as a \texttt{Message} keyed by \texttt{group\_id} and broadcasts to connected group members in real time; if the socket is unavailable the client falls back to authenticated HTTP polling of the group's messages.}
\label{fig-seq-messaging}
\end{figure}

% auto-generated figure: seq-ingest (source: fix)
\begin{figure}[H]
\centering
{\footnotesize
\begin{sequencediagram}
  \newthread{sch}{:Scheduler}
  \newinst[1]{svc}{:training-match}
  \newinst[1]{api}{football-data.org}
  \newinst[1]{db}{:matches DB}
  \newinst[1]{fe}{:Frontend}
  \mess{sch}{every 30 min / sync}{svc}
  \begin{call}{svc}{GET matches + standings}{api}{collections}
  \end{call}
  \begin{sdblock}{loop}{filter to club, de-dup by external id}
    \begin{call}{svc}{upsert match}{db}{ok}
    \end{call}
  \end{sdblock}
  \mess{svc}{standings read-through}{svc}
  \begin{call}{fe}{GET /matches}{svc}{fixtures + table}
  \end{call}
\end{sequencediagram}
}
\caption{Sequence diagram of scheduled external fixture ingestion: the training-match service pulls competition-scoped matches and standings from football-data.org, filters and de-duplicates fixtures before persisting them, reads standings through without storing, and lets the front-end consume the data from the platform's own model.}
\label{fig-seq-ingest}
\end{figure}

% auto-generated figure: seq-notify (source: fix)
\begin{figure}[H]
\centering
{\footnotesize
\begin{sequencediagram}
  \newthread{ds}{:domain service}
  \newinst[1]{kf}{:Kafka}
  \newinst[1]{nm}{:notification-mail}
  \newinst[1]{db}{:DB}
  \newinst[1]{fe}{:Recipient}
  \mess{ds}{publish event / outbox}{kf}
  \mess{kf}{deliver async}{nm}
  \begin{call}{nm}{persist notification}{db}{saved}
  \end{call}
  \mess{fe}{GET /notifications}{nm}
  \mess{nm}{unread list + bell badge}{fe}
\end{sequencediagram}
}
\caption{Sequence diagram of event-driven notification delivery: the training-match service publishes a domain event to Kafka and returns immediately, while the notification-mail service asynchronously consumes the event, persists a notification, and later surfaces it to the recipient's top-bar bell.}
\label{fig-seq-notify}
\end{figure}

% auto-generated figure: seq-password (source: wf)
\begin{figure}[H]
\centering
{\footnotesize
\begin{sequencediagram}
  \newthread{u}{:User}
  \newinst[1]{fe}{:Frontend}
  \newinst[1]{gw}{:Gateway}
  \newinst[1]{um}{:user-management}
  \newinst[1]{kc}{:Keycloak}

  \begin{call}{u}{submit current + new password}{fe}{form reset}
    \begin{call}{fe}{PUT /users/\{id\}/password}{gw}{result}
      \begin{call}{gw}{route to user-management}{um}{result}
        \begin{call}{um}{password-grant token (current pwd)}{kc}{token / error}
        \end{call}
        \begin{sdblock}{alt}{[grant fails: 401 invalid\_grant]}
          \mess{um}{Current password is incorrect}{gw}
        \end{sdblock}
        \begin{sdblock}{}{[grant ok: current password verified]}
          \begin{call}{um}{Admin API reset-password (new pwd)}{kc}{204 No Content}
          \end{call}
          \mess{um}{200 OK password changed}{gw}
        \end{sdblock}
      \end{call}
    \end{call}
  \end{call}
\end{sequencediagram}
}
\caption{Sequence diagram of the password-change flow, in which the user-management service first verifies the current password via an OAuth~2.0 password-grant exchange against Keycloak and only on success sets the new password through the Keycloak Admin API.}
\label{fig-seq-password}
\end{figure}


%==============================================================================
\section{State and Lifecycle Models}
\label{sec:design-state}

Finally, three entities have a life cycle rich enough to deserve an explicit
state machine. Figure~\ref{fig-state-match} models the broadcast-style phase
machine of a live match, including the extra-time and penalty branch reserved
for unresolved knockout ties. Figure~\ref{fig-state-injury} models the injury
life cycle, whose stage is read directly from the injury's status and which
returns a player to availability only on a passing fitness test.
Figure~\ref{fig-state-competition} models the competition life cycle, including
the Champions League specialisation in which the knockout bracket is locked
until the league phase is complete and a finished competition becomes read-only.

% auto-generated figure: state-match (source: fix)
\begin{figure}[H]
\centering
\resizebox{0.95\textwidth}{!}{%
\begin{tikzpicture}[node distance=14mm and 16mm,
    every edge/.style={draw=dLine,thick,-{Stealth[length=2.4mm]}}]
  \node (start) {};
  \node[stbox, right=8mm of start] (sch) {Scheduled};
  \node[stbox, right=of sch] (h1) {First\\Half};
  \node[stbox, right=of h1]  (ht) {Half\\Time};
  \node[stbox, right=of ht]  (h2) {Second\\Half};
  \node[stbox, right=of h2]  (ft) {Full\\Time};
  \node[stbox, below=14mm of ft]  (et)  {Extra\\Time};
  \node[stbox, below=14mm of et]  (pen) {Penalties};
  \node[stbox, below=14mm of pen, fill=dFill2, draw=dLine2] (fin) {Finished};
  \draw (start) edge (sch);
  \draw (sch) edge node[lbl,above]{kick-off} (h1);
  \draw (h1)  edge node[lbl,above]{45$'$} (ht);
  \draw (ht)  edge node[lbl,above]{resume} (h2);
  \draw (h2)  edge node[lbl,above]{90$'$} (ft);
  \draw (ft)  edge node[lbl,left]{knockout \& level} (et);
  \draw (et)  edge node[lbl,left]{still level} (pen);
  \draw (pen) edge node[lbl,left]{shoot-out} (fin);
  \draw (ft)  to[out=0,in=0,looseness=1.5] node[lbl,right]{league / draw OK} (fin);
\end{tikzpicture}%
}
\caption{State machine for live-match phase progression, where a scheduled match advances through the regulation halves to full time and, for an unresolved knockout tie only, continues into extra time and a penalty shoot-out before finishing.}
\label{fig-state-match}
\end{figure}

% auto-generated figure: state-injury (source: fix)
\begin{figure}[H]
\centering
\resizebox{\textwidth}{!}{%
\begin{tikzpicture}[every edge/.style={draw=dLine,thick,-{Stealth[length=2.4mm]}}]
  \node (start) {};
  \node[stbox, right=10mm of start] (act)  {Active\\\scriptsize status=active};
  \node[stbox, right=16mm of act]   (diag) {Diagnosis};
  \node[stbox, right=16mm of diag]  (treat){Treatment\\\scriptsize $\rightarrow$ recovering};
  \node[stbox, right=16mm of treat] (rehab){Rehabili-\\tation};
  \node[stbox, right=16mm of rehab] (prog) {Recovery\\Program};
  \node[stbox, right=16mm of prog]  (fit)  {Fitness\\Test};
  \node[stbox, right=16mm of fit, fill=dFill2, draw=dLine2] (recd) {Recovered\\\scriptsize available};
  \draw (start) edge node[lbl,above]{injury} (act);
  \draw (act)   edge node[lbl,above]{diagnose} (diag);
  \draw (diag)  edge node[lbl,above]{treat} (treat);
  \draw (treat) edge node[lbl,above]{rehab} (rehab);
  \draw (rehab) edge node[lbl,above]{programme} (prog);
  \draw (prog)  edge node[lbl,above]{test} (fit);
  \draw (fit)   edge node[lbl,above]{pass} (recd);
  \draw (fit)   edge[loop below] node[lbl,below]{fail} (fit);
\end{tikzpicture}%
}
\caption{State machine for the injury lifecycle and recovery journey: the Injury.status field drives Active to Recovering to Recovered, with the staged journey (Diagnosis, Treatment, Rehabilitation, Recovery Program, Fitness Test) walked along the transitions, a failed fitness test looping back and a passed test restoring the player to Available.}
\label{fig-state-injury}
\end{figure}

% auto-generated figure: state-competition (source: wf)
\begin{figure}[H]
\centering
\resizebox{\textwidth}{!}{%
\begin{tikzpicture}[node distance=26mm, on grid,
    every edge/.style={flow,draw}]
  % --- main competition lifecycle ---
  \node[st,initial,initial text=create] (draft) {Draft};
  \node[st,right=34mm of draft]         (ongo)  {Ongoing\\(editable)};
  \node[st,accepting,right=40mm of ongo](comp)  {Completed\\(read-only)};

  \draw[flow] (draft) edge node[lbl,above]{publish} (ongo);
  \draw[flow] (ongo)  edge node[lbl,above,align=center]{all fixtures\\played} (comp);
  \draw[flow] (ongo)  edge[loop above] node[lbl]{record result} (ongo);

  % --- Champions League sub-states (knockout gated by league phase) ---
  \node[st,below=24mm of draft] (league) {League\\Phase};
  \node[st,right=44mm of league] (knock) {Knockout\\(locked)};

  \draw[flow] (league) edge node[lbl,above,align=center]{league phase complete\\$\rightarrow$ unlock knockout} (knock);
  \draw[flow] (knock)  edge[bend left=12] node[lbl,below,align=center]{all ties played} (comp);

  % link the CL specialisation to the generic Ongoing state
  \draw[flowd] (ongo) edge[bend right=14] node[lbl,right,align=center]{Champions\\League} (league);

  \node[font=\scriptsize\itshape,below=10mm of league,xshift=22mm,align=center]
    {Champions League specialisation of \emph{Ongoing}:\\the knockout bracket is locked until the league phase is complete.};
\end{tikzpicture}}
\caption{State machine for the competition lifecycle: a competition moves from Draft through an editable Ongoing phase to a read-only Completed state once every fixture is played, with the Champions League specialisation in which the knockout bracket stays locked until the league phase is complete.}
\label{fig-state-competition}
\end{figure}

